% 12-functions.tex — Clause 12: Functions and procedures
\chapter{Functions and procedures}
\label{sec:12-functions}
\index{function}
\index{procedure}

\section{The func/proc distinction}
\label{sec:fn.distinction}
\index{func vs proc}

\pnum
Lain distinguishes between pure functions (\lain{func}) and side-effecting
procedures (\lain{proc}).  This distinction is enforced statically and
reflected in the type system.

\begin{longtable}{@{}lll@{}}
\toprule
\textbf{Property} & \textbf{\lain{func}} & \textbf{\lain{proc}} \\
\midrule
\endfirsthead
\toprule
\textbf{Property} & \textbf{\lain{func}} & \textbf{\lain{proc}} \\
\midrule
\endhead
\bottomrule
\endlastfoot
Side effects                      & forbidden          & allowed \\
Unbounded \lain{while}            & forbidden          & allowed \\
Bounded \lain{while decreasing}   & allowed            & allowed \\
\lain{for} loops                  & allowed            & allowed \\
Direct/indirect recursion         & forbidden          & allowed \\
Calling \lain{proc}               & forbidden          & allowed \\
Calling \lain{func}               & allowed            & allowed \\
Global mutable state              & no access          & full access \\
Termination guarantee             & yes                & no \\
\end{longtable}

\pnum
A \lain{func} is referentially transparent: calling it with the same
arguments always produces the same result.  Its evaluation has no
observable effect on any state outside its own stack frame.

\pnum
Local mutation through a \lain{var T} parameter (exclusive mutable
borrow) is permitted in a \lain{func}. The borrow checker guarantees
no aliasing, so the mutation is observable only through the
parameter binding itself; the function remains deterministic and
referentially transparent (same input borrow $\Rightarrow$ same
mutation pattern).

\begin{example}
\begin{laincode}
type Counter { val i32 }

func increment(var c Counter) {  // legal: mutation is local
    c.val = c.val + 1
}
\end{laincode}
\end{example}

\pnum
The compiler emits \diagcode{W130} when a \lain{proc} has no
observable side effect — no calls to other \lain{proc} or
\lain{extern proc}, no \lain{panic}, no mutable global access, no
unbounded \lain{while}, no self-recursion — to suggest that it could
be declared \lain{func}. This is a non-blocking warning; the
programmer may keep it as \lain{proc} when future evolution will
introduce effects.

\section{Purity enforcement}
\label{sec:fn.purity}
\index{purity enforcement}

\begin{constraint}
  A \lain{func} that calls a \lain{proc} or \lain{extern proc} is
  \illformed{}~(\diagcode{E011}).
\end{constraint}

\begin{constraint}
  A \lain{func} that contains an unbounded \lain{while} loop (one without
  a \lain{decreasing} measure) is \illformed{}~(\diagcode{E011}).
\end{constraint}

\begin{constraint}
  A \lain{func} that calls itself directly, or that participates in a
  mutual-recursion cycle, is \illformed{}~(\diagcode{E011}).  The
  implementation shall detect mutual recursion by performing a depth-first
  search on the call graph.
\end{constraint}

\begin{note}
  Generic functions (\lain{func} with \lain{type} parameters) are
  excluded from mutual-recursion detection until they are monomorphized.
\end{note}

\section{Termination analysis}
\label{sec:fn.termination}
\index{termination analysis}

\pnum
A \lain{func} is guaranteed to terminate because:
\begin{enumerate}
  \item all \lain{while} loops carry a verified termination measure;
  \item \lain{for} loops iterate over finite ranges;
  \item recursion is forbidden.
\end{enumerate}

\pnum
No termination analysis is performed for \lain{proc}s.

\section{Universal Function Call Syntax (UFCS)}
\label{sec:fn.ufcs}
\index{UFCS}

\pnum
An expression of the form \lain{x.f(args)} is desugared to \lain{f(x, args)}
before name resolution, provided \lain{f} is not a declared field of
\lain{x}'s type.  This is called Universal Function Call Syntax (UFCS).

\pnum
\textbf{UFCS ownership}: the ownership mode of the first argument in the
desugared call is determined by the declared mode of the first parameter of
\lain{f}:

\begin{itemize}
  \item If the first parameter is \lain{var}, the receiver is automatically
        mutably borrowed: \lain{x.f()} desugars to \lain{f(var x)}.
  \item If the first parameter is \lain{mov}, the receiver is automatically
        moved: \lain{x.f()} desugars to \lain{f(mov x)}.
  \item If the first parameter is shared (default), the receiver is shared:
        \lain{x.f()} desugars to \lain{f(x)}.
\end{itemize}

\pnum
\textbf{UFCS resolution algorithm}:
\begin{enumerate}
  \item Check if \lain{f} is a declared field of \lain{x}'s type.  If so,
        this is a field access, not UFCS.
  \item Search for a function named \lain{f} whose first parameter type is
        compatible with the type of \lain{x}.
  \item If found, desugar the call.
  \item If no matching function is found, the call is unresolved and the
        program is \illformed{} (an unresolved receiver or callee is reported;
        a bare undeclared identifier is \diagcode{E106},
        \S\,\ref{sec:basics.resolution}).
\end{enumerate}

\begin{example}
\begin{laincode}
func is_empty(s u8[]) bool { return s.len == 0 }

var buf = "hello"
var empty = buf.is_empty()   // desugars to: is_empty(buf)
\end{laincode}
\end{example}

\section{Calling conventions}
\label{sec:fn.abi}
\index{calling convention}

\pnum
The calling convention for Lain functions corresponds to the C99 calling
convention of the target platform, as determined by the C compiler used for
the translation output.  The parameter passing rules are:

\begin{longtable}{@{}lll@{}}
\toprule
\textbf{Lain parameter} & \textbf{C type} & \textbf{Semantics} \\
\midrule
\endfirsthead
\toprule
\textbf{Lain parameter} & \textbf{C type} & \textbf{Semantics} \\
\midrule
\endhead
\bottomrule
\endlastfoot
shared struct \lain{T} & \texttt{const T*} & pointer to read-only copy \\
\lain{var} struct \lain{T} & \texttt{T*}       & pointer, caller's copy \\
\lain{mov} struct \lain{T} & \texttt{T}        & by value (bitwise copy) \\
primitive any mode      & \texttt{T}        & by value \\
\end{longtable}

\section{Return values}
\label{sec:fn.return}
\index{return value}

\pnum
A return statement has the form \lain{return <expr>}.  The mode of the
returned value is determined by the function's declared return type and
follows the same rules as call-site argument passing.

\pnum
A \lain{return} may be prefixed with \lain{mov} or \lain{var}:
\lain{return mov x} makes the ownership transfer explicit, and
\lain{return var x} returns a mutable borrow (the pattern for a \lain{func}
whose declared return type is \lain{var T}).  When the return type is owned,
a bare \lain{return expr} already moves \lain{expr}, so \lain{mov} is optional
there.

\begin{note}
  A \lain{return var x} / \lain{return mov x} that would escape a borrow of a
  \emph{local} variable is a dangling reference (\diagcode{E010},
  \S\,\ref{sec:stmt.return}).  This supersedes an earlier design (Q-009) in
  which explicit \lain{mov}/\lain{var} in return position was rejected.
\end{note}


\section{The \texttt{panic} builtin}
\label{sec:fn.panic}
\index{panic}

\pnum
\lain{panic(msg)} is a builtin expression that terminates program
execution immediately.  Its type is \textit{Never} (the bottom type:
a subtype of every type).  Because it does not return, \lain{panic}
may appear in any expression position---including as the body of a
match arm whose declared type is otherwise unrelated.

\pnum
\lain{panic} may be called from both \lain{func} and \lain{proc}
contexts.  The call is deterministic (same input yields the same
termination), so it does not violate the purity of an enclosing
\lain{func}.

\begin{constraint}
  \textbf{Q-014, S16.}
  When \lain{panic} is executed, the program terminates immediately
  via the platform abort primitive.  No \lain{defer} block is
  executed.  No stack unwinding occurs.
\end{constraint}

\begin{note}
  The reference implementation emits, for \lain{panic(msg)}, the
  equivalent of
  \texttt{fprintf(stderr, "panic: \%.*s\textbackslash{}n", ...) ; abort()}.
  The message is written to the standard error stream before abort.
\end{note}

\begin{example}
\begin{laincode}
func head(s u8[]) u8 {
    if s.len == 0 {
        panic("empty slice")    // Never; satisfies the return type
    }
    return s[0]
}
\end{laincode}
\end{example}

\section{Function attributes}
\label{sec:fn.attributes}
\index{attribute}

\pnum
A function declaration may be prefixed with one or more attributes
(\S\,\ref{sec:decl.attributes}).  The attributes applicable to
functions in version~1.0 are:

\begin{itemize}
  \item \tok{[fast\_math]} --- enable fast floating-point optimizations
        (FMA fusion, contraction) for the body of this function.  This
        is an opt-out from the determinism guarantee (P4) and applies
        only locally; functions without \tok{[fast\_math]} remain
        bit-deterministic within their target platform.
  \item \tok{[private]} --- module-private visibility
        (\S\,\ref{sec:mod.visibility}).
\end{itemize}

\begin{note}
  The reference implementation translates \tok{[fast\_math]} on a
  function~$f$ to a pair of pragmas surrounding the body of $f$ in the
  emitted C source, enabling \texttt{fp-contract=fast} for $f$ only.
  The default (i.e., absence of \tok{[fast\_math]}) emits
  \texttt{-ffp-contract=off} flags, providing deterministic floating
  point within a single target platform and C compiler.
\end{note}
